\begin{table}
\centering
\begin{tabular}{lccccc}
\toprule
& conv1 & conv2 & conv3 & conv4 & conv5 \\
\midrule
Random \citep{jigsaw}            & 48.5 & 41.0 & 34.8 & 27.1 & 12.0 \\
\citet{xiaolong}                 & 51.8 & 46.9 & 42.8 & 38.8 & 29.8 \\
\citet{carl}                     & 53.1 & 47.6 & 48.7 & \bf{45.6} & 30.4 \\
\gray{\citet{jigsaw}}*           & \gray{57.1} & \gray{56.0} & \gray{52.4} & \gray{48.3} & \gray{38.1} \\
\midrule
BiGAN (ours)                     & \bf{56.2} & \bf{54.4} & \bf{49.4} & 43.9 & 33.3 \\
BiGAN, $112\times112$ $E$ (ours) & 55.3 & 53.2 & 49.3 & 44.4 & \bf{34.8} \\
\bottomrule
\end{tabular}
\caption{
Classification accuracy (\%) for the ImageNet LSVRC~\citep{imagenet} validation set with various portions of the network frozen, or reinitialized and trained from scratch, following the evaluation from~\citet{jigsaw}.
In, e.g., the \textit{conv3} column, the first three layers -- conv1 through conv3 --
are transferred and frozen, and the last layers --
conv4, conv5, and fully connected layers --
are reinitialized and trained fully supervised for ImageNet classification.
BiGAN is competitive with these contemporary visual feature learning methods, despite its generality.
(*Results from \citet{jigsaw} are not directly comparable to those of the other methods
as a different base convnet architecture with larger intermediate feature maps is used.)
}
\lbltbl{imagenetclass}
\end{table}
